\section*{Chapitre \ref{ch:ml:interpretability-and-model-governance} --- Interprétabilité et gouvernance des modèles}

\begin{solution}{exo:ml:interpretability-and-model-governance:1}
$\mathrm{PD}(v) = v + \E[x_2^2] = v + 1/3$~: l'effet réel plus une constante. Si $x_2 = x_1$ et que le modèle est $f$ lui-même,
la \omterm{def:ml:interpretability-and-model-governance:pd}{dépendance partielle} moyenne encore $x_2^2$ sur la marginale de $x_2$ et renvoie $v + 1/3$, alors que le long des données la
prédiction est $v + v^2$~: la \omterm{def:ml:interpretability-and-model-governance:pd}{dépendance partielle} décrit la formule, et non ce que fait le modèle sur les données, et ne
attribue aucune part de la courbure à la seule variable qui bouge.
\end{solution}

\begin{solution}{exo:ml:interpretability-and-model-governance:2}
$0.1\times250 = 25$ attendus. La probabilité d'obtenir 27 ou plus est de 0.37~: rien d'alarmant.
\end{solution}

\begin{solution}{exo:ml:interpretability-and-model-governance:3}
Les règles de l'arbre décrivent la part du comportement du modèle qu'elles reproduisent, soit moins de la moitié de la variance ici~; rien ne peut être
conclu sur le reste, qui est ici toute l'interaction. Les règles d'un substitut ne sont une explication que là où sa fidélité
est élevée, et doivent être vérifiées sur les tranches pertinentes.
\end{solution}

\begin{solution}{exo:ml:interpretability-and-model-governance:4}
Son effet vaut $2x_3x_4$, avec $x_4$ symétrique autour de zéro~: pour chaque observation où la hausse de $x_3$ fait monter la prédiction, il
en existe une où elle la fait baisser, et la pente moyenne est nulle. Les pentes des courbes ICE suivent $x_4$.
\end{solution}

\begin{solution}{exo:ml:interpretability-and-model-governance:5}
Les cinq variables faibles ont des effets réels égaux~; leurs importances estimées diffèrent par le bruit d'échantillonnage, ce qui les réordonne dans
chaque rééchantillon. Seules les différences d'importance plus grandes que ce bruit sont stables.
\end{solution}

\begin{solution}{exo:ml:interpretability-and-model-governance:6}
Une part d'attribution SHAP issue d'un seul ajustement, avec des variables corrélées et sans contrôle de stabilité, ne dit pas que le modèle est un
modèle de révisions~: les variables corrélées se partagent le mérite de façon arbitraire, la part peut changer au prochain réentraînement, et l'attribution n'est pas
une causalité. Montrer l'ALE et l'ICE de la variable, l'attribution d'un réentraînement à l'autre, et la performance du modèle sans elle.
\end{solution}

\begin{solution}{exo:ml:interpretability-and-model-governance:7}
Il semble atteindre la perfection~: les cinq premières sont les variables 0 à 4 à chaque réentraînement et la corrélation de Kendall vaut 1.
L'explication est que rien n'a été appris~: avec 400 observations, un \omterm{def:ml:trees-and-boosting:bagging}{bagging} à 70~\% et le minimum de 200 observations par
feuille fixé par l'établissement, aucun arbre ne peut se scinder~; chaque \omterm{def:ml:feature-engineering-selection-and-importance:shapley}{valeur SHAP} est nulle, et le classement retombe sur l'ordre des colonnes. Un score de stabilité doit être lu
avec les importances qu'il classe.
\end{solution}

\begin{solution}{exo:ml:interpretability-and-model-governance:8}
L'ALE de $j$ cumule $\E[f(z_{k}, x_{-j}) - f(z_{k-1}, x_{-j})\mid x_j\in(z_{k-1}, z_k]]$~; pour un $f$ additif, les autres
composantes s'annulent dans chaque différence, ce qui laisse $f_j(z_k) - f_j(z_{k-1})$, et la somme se télescope en $f_j$ à une constante près.
La \omterm{def:ml:interpretability-and-model-governance:pd}{dépendance partielle} vaut $f_j(v) + \sum_{k\ne j}\E f_k(x_k)$, soit aussi $f_j$ à une constante près. Les deux sont justes pour la vraie
fonction~; la \omterm{def:ml:interpretability-and-model-governance:pd}{dépendance partielle} se trompe pour un modèle ajusté, qui n'est contraint que là où sont les données et peut ne pas être
additif hors du support des données, là où la \omterm{def:ml:interpretability-and-model-governance:pd}{dépendance partielle} l'évalue.
\end{solution}

\begin{solution}{pb:ml:interpretability-and-model-governance:1}
\textbf{Partie I.}
\begin{enumerate}
\item Aux corrélations 0, 0.92 et 0.99~: PD 0.027, 0.134, 0.229~; ALE 0.028, 0.044, 0.110.
\item Le modèle ajusté ne peut pas séparer des variables presque identiques~; l'ALE explique le modèle, qui mêle leurs effets.
\item Les pentes se dispersent avec un écart-type de 0.91 et une corrélation de 0.98 avec $x_4$~: l'interaction.
\item La partie additive en $x_1$ et $x_2$ (fidélité 0.48)~; aucune part de l'interaction.
\end{enumerate}
\textbf{Partie II.}
\begin{enumerate}[resume]
\item La prédiction (0.82) provient surtout de l'interaction~; $x_1$ devrait passer de 0.46 à $-0.44$ pour l'inverser.
\item À $R^2$ 0.66~: le même ensemble dans 22~\% des réentraînements, Kendall 0.89~; à $R^2$ 0.05~: jamais, Kendall 0.78, avec des variables nulles
parfois parmi les cinq premières.
\item Importances issues de plusieurs réentraînements, leur dispersion, les variables stables en tête, et une mise en garde sur les groupes corrélés.
\item Avec les \omterm{def:ml:feature-engineering-selection-and-importance:shapley}{valeurs SHAP} du modèle déployé pour cette action le lundi et le mardi, les entrées qui ont changé, une \omterm{def:ml:interpretability-and-model-governance:cf}{explication contrefactuelle},
et l'ALE de ces entrées montrant que la réponse est cohérente.
\end{enumerate}
\textbf{Partie III.}
\begin{enumerate}[resume]
\item Quinze rubriques, du nom aux approbations~; le suivi et les approbations sont absents.
\item Performance hors échantillon ($R^2$ 0.860), le challenger (régression ridge 0.349), parité d'implémentation (exacte sur 200 lignes) et
l'analyse des résultats (27 sur 250, queue 0.37).
\item Le modèle en production doit être celui qui a été validé~; les exports, les compilateurs (chapitre~26) et les réentraînements rompent cette correspondance silencieusement.
\item Que les superviseurs traitent les modèles appris comme des modèles, selon les mêmes principes, quelle que soit la technologie.
\end{enumerate}
\textbf{Partie IV.}
\begin{enumerate}[resume]
\item \emph{Expliquer le résultat au comité des risques.} L'erreur de la \omterm{def:ml:interpretability-and-model-governance:pd}{dépendance partielle} par rapport à l'effet réel vaut 0.134 et 0.229 pour des
corrélations de 0.92 et 0.99, contre 0.044 et 0.110 pour l'ALE~; les cinq premières variables SHAP forment le même ensemble dans 22~\% des dix réentraînements
à $R^2$ 0.66 (Kendall 0.89) et dans aucun à $R^2$ 0.05 (Kendall 0.78).
\item Non ici, où l'interaction porte la moitié du signal et où le challenger la perd~; oui là où un modèle linéaire reste
à la portée du bruit du modèle boosté.
\item Un niveau élevé s'il pilote des positions à grande échelle~: significatif, complexe et incertain.
\item Une équipe indépendante reconstruisant le modèle à partir de la fiche, un challenger, des entrées de stress hors de la plage d'entraînement, et
des explications vérifiées au regard du bon sens économique.
\item Performance, stabilité, explications et parité~; la version de la fiche et les champs de données.
\item Au niveau des variables qui alimentent le plongement, ou au moyen de sondes reliant les directions du plongement à des variables connues.
\item Les secrets (identifiants, données clients) et les affirmations sans preuve.
\item Un modèle du modèle, avec une erreur qui lui est propre.
\end{enumerate}
\end{solution}

\begin{solution}{iq:ml:interpretability-and-model-governance:1}
La \omterm{def:ml:interpretability-and-model-governance:pd}{dépendance partielle} moyenne les prédictions en fixant la variable à chaque valeur pour chaque observation, en parcourant des combinaisons
absentes des données~; l'ALE moyenne les différences locales dans des intervalles, pour des observations réellement présentes. Avec des
variables corrélées, la \omterm{def:ml:interpretability-and-model-governance:pd}{dépendance partielle} peut être très fausse~; l'ALE reste sur les données.
\iqlookfor{les deux définitions et le problème hors variété des données.}
\end{solution}

\begin{solution}{iq:ml:interpretability-and-model-governance:2}
Souvent instable en dehors des toutes premières variables, surtout avec des signaux faibles et des variables corrélées. Mesurer en réentraînant sur
des rééchantillons ou des périodes successives, et publier l'accord sur l'ensemble de tête et les corrélations de rang, à côté des importances elles-mêmes.
\iqlookfor{instabilité de la queue du classement et protocole de mesure.}
\end{solution}

\begin{solution}{iq:ml:interpretability-and-model-governance:3}
Objet et usage, cible et horizon, données et période, variables, procédure d'entraînement, résultats de validation avec un challenger,
performance par régime, explications, limites, suivi, responsables et approbations, version.
\iqlookfor{objet, données, évaluation, limites et responsabilité.}
\end{solution}

\begin{solution}{iq:ml:interpretability-and-model-governance:4}
La variation de la prévision et les entrées qui l'ont produite (SHAP sur le modèle déployé), si la réponse est cohérente (ALE/ICE),
ce qui l'aurait inversée, et comment la règle de position a transformé la prévision en transaction.
\iqlookfor{attribution locale, cohérence de la réponse et lien avec la position.}
\end{solution}

\begin{solution}{iq:ml:interpretability-and-model-governance:5}
Performance hors période face à un challenger simple, contrôles de fuite de données, stabilité d'un réentraînement à l'autre, explications vérifiées quant au signe
et à l'ampleur, parité d'implémentation du modèle exporté, analyse des résultats, une \omterm{def:ml:interpretability-and-model-governance:card}{fiche de modèle}, et suivi avant la mise en production.
\iqlookfor{challenger, fuite de données, stabilité, parité et suivi.}
\end{solution}

\begin{solution}{iq:ml:interpretability-and-model-governance:6}
Lorsque le modèle interprétable reste dans le bruit de la boîte noire, lorsque les décisions doivent être justifiées cas par cas, et lorsque le
coût d'une défaillance non expliquée est élevé~; les signaux faibles du trading rendent souvent la première condition vraie.
\iqlookfor{parité de performance et coût de l'opacité.}
\end{solution}
